% Source: 02 - Tue Sep 29 2026.pdf, page 15. Content remains in source order.
\sourcepage{02}{15}
\section{Polynomial-time verification of Hamiltonian cycles}

Given $G=(V,E)$ undirected, a simple cycle is a sequence:
\[
\langle v_1,v_2,\ldots,v_k\rangle,\quad v_i\in V,\quad v_1=v_k,
\]
\[
v_i\neq v_j\quad1\leq i<j<k,\qquad k\geq1,\qquad
\{v_i,v_{i+1}\}\in E\quad1\leq i<k.
\]

HAMILTON problem is a decision problem that asks whether a given undirected graph contains a simple cycle that visits all nodes in the graph exactly once (a Hamiltonian cycle):
\[
\begin{aligned}
\text{I: }&\langle G=(V,E)\rangle,\quad G\text{ undirected},\\
\text{D: }&\text{is there in }G\text{ a simple cycle containing all nodes in }V?
\end{aligned}
\]
Example:
\sourcefigure[0.65]{assets/02/page-015-hamilton-cycle.png}

\subsection{Solution}
Deciding whether a graph contains a Hamiltonian cycle is a difficult problem (no polynomial algorithms are known).
\subsection{Verification}
Verifying that a candidate vertex sequence is indeed a Hamiltonian cycle is easy (there is a polynomial time algorithm that checks the correctness of a candidate solution). Given a candidate sequence $\langle v_1,v_2,\ldots,v_k\rangle$, we can check in polynomial time whether it is a simple cycle and whether it visits all nodes in the graph exactly once. If both conditions are satisfied, we can conclude that the graph contains a Hamiltonian cycle.
